% !TeX program = xelatex
% 시연 영상용 개요·성과 슬라이드 (16:9). build_slides.sh가 본문 그림을 추출한 뒤 빌드하고 PNG로 내보낸다.
\documentclass[aspectratio=169,11pt]{beamer}
\usepackage{kotex}
\setmainfont{Noto Sans CJK KR}
\setsansfont{Noto Sans CJK KR}
\newfontfamily\figsans{Noto Sans CJK KR}
\usepackage{booktabs,tabularx,array,multirow}
\usepackage{tikz}
\usetikzlibrary{arrows.meta,positioning,shadows,calc,backgrounds}
\usepackage{pgfplots}
\pgfplotsset{compat=1.16}
\pgfkeys{/pgf/number format/assume math mode=true}  % 그래프 숫자를 본문 글꼴로
\newcolumntype{Y}{>{\raggedright\arraybackslash}X}
\renewcommand{\arraystretch}{1.25}

\definecolor{ink}{HTML}{1E2024}
\definecolor{muted}{HTML}{6E7278}
\definecolor{bg}{HTML}{FAFAF8}
\definecolor{afs}{HTML}{234C70}
\definecolor{rnd}{HTML}{B0BEC5}
\definecolor{passc}{HTML}{228B54}
\definecolor{failc}{HTML}{C8322D}
\setbeamercolor{background canvas}{bg=bg}
\setbeamercolor{normal text}{fg=ink}
\setbeamercolor{frametitle}{fg=ink,bg=bg}
\setbeamercolor{structure}{fg=afs}
\setbeamerfont{frametitle}{size=\LARGE,series=\bfseries}
\setbeamertemplate{navigation symbols}{}
\setbeamertemplate{itemize items}{\color{afs}\textbullet}
\setbeamertemplate{frametitle}{\vspace*{5mm}\insertframetitle\par
  \ifx\insertframesubtitle\empty\else{\normalsize\color{muted}\mdseries\insertframesubtitle\par}\fi}
\setbeamertemplate{footline}{\hfill{\scriptsize\color{muted}\insertframenumber\hspace*{6mm}}\vspace*{3mm}}
\setbeamersize{text margin left=10mm,text margin right=10mm}
\newcommand{\bignum}[2]{{\color{#2}\fontsize{30}{32}\selectfont\bfseries #1}}

\begin{document}

% ------------------------------------------------------------------ 1 표지
{\setbeamertemplate{footline}{}
\begin{frame}
\vspace*{10mm}
{\fontsize{22}{32}\selectfont\bfseries AI 기반 Physical AI 작업 수행 검증 및\\회귀 테스트 자동화 기술\par}
\vspace{6mm}
{\large\color{muted} 행동 피드백 기반 능동 실패 탐색(AFS) 시연\par}
\vspace{14mm}
{\normalsize 실험 1 \; 로봇팔 pick-and-place: Franka Panda 로봇팔 + GR00T N1.7 VLA 정책\\[1mm] 실험 2 \; 휴머노이드 이동: Unitree G1 휴머노이드 + LLM Agent\par}
\end{frame}}

% ------------------------------------------------------------------ 2 배경과 문제
\begin{frame}{배경과 문제}{같은 명령도 환경이 조금 바뀌면 실패한다}
\vspace{4mm}
\begin{columns}[T]
\begin{column}{0.5\textwidth}
\textbf{Physical AI의 특성}
\begin{itemize}
\item 로봇팔, 자율 이동 로봇, VLA 모델이 현장에 도입
\item 실제 공간에서 행동 → 소프트웨어 정확도만으로 품질 보장 불가
\item 물체 위치, 장애물, 작업자 접근에 따라 결과가 달라짐
\end{itemize}
\end{column}
\begin{column}{0.46\textwidth}
\textbf{현재 검증 방식의 한계}
\begin{itemize}
\item 예상 가능한 케이스를 사람이 수작업 정의
\item 실제 장비 반복은 시간·충돌 위험
\item 드문 실패, PASS/FAIL 경계를 체계적으로 찾기 어려움
\item 현장 배치가 바뀔 때마다 테스트 재작성
\end{itemize}
\end{column}
\end{columns}
\vspace{8mm}
\begin{center}
\colorbox{afs!10}{\parbox{0.9\textwidth}{\centering\large
\textbf{목표}: 실패할 가능성이 높은 테스트 환경을 자동으로 찾아내고, 찾은 실패를 회귀 테스트로 재사용}}
\end{center}
\end{frame}

% ------------------------------------------------------------------ 3 목표 지표
\begin{frame}{과제 평가 지표}{같은 시도 횟수에서 Random 테스트와 비교}
\vspace{6mm}
\begin{tabularx}{\textwidth}{@{}p{40mm}Yp{24mm}@{}}
\toprule
지표 & 정의 & 목표 \\
\midrule
실패 발견율 (FDR) & 유효 테스트 중 작업 실패를 발견한 비율 & \textbf{30\% 이상} \\
Random 테스트 대비 향상 & 같은 시도 횟수에서 Random 테스트 대비 실패 발견 향상률 & \textbf{20\% 이상} \\
실패 유형 다양성 & 사전 정의한 주요 실패 유형 6종 중 발견 수 & \textbf{4종 이상} \\
\bottomrule
\end{tabularx}
\vspace{8mm}

{\color{muted} 실패 유형 6종: 충돌 · 도달 불가 · 장애물 간섭 · 목표 위치 점유 · 작업자 안전 위험 · 인식 오류}
\end{frame}

% ------------------------------------------------------------------ 4 시스템 구성
\begin{frame}{시스템 구성}{탐색(AFS) · 실행 · 판정을 분리하고, 테스트 대상 로봇은 블랙박스로 실행}
\centering
\resizebox{!}{0.84\textheight}{\input{build/fig_arch.tex}}
\end{frame}

% ------------------------------------------------------------------ 5 적용 기술
\begin{frame}{주요 적용 기술}
\vspace{4mm}
\small
\begin{tabularx}{\textwidth}{@{}p{34mm}Y@{}}
\toprule
구성 요소 & 적용 기술 \\
\midrule
장면 표현 & 장면 파라미터 벡터, 3D 장면 그래프, 점유 격자 이동 지도 \\
물리 시뮬레이션 & MuJoCo, MJCF 자동 생성, robosuite·LIBERO \\
테스트 대상 & Panda 로봇팔 + GR00T N1.7 VLA 정책 \quad / \quad G1 휴머노이드 + LLM Agent + GR00T 보행 제어 \\
능동 실패 탐색 & GPT-6 Luna Pro / Luna (OpenRouter), 구조화된 입출력, 규칙 기반 후보 검증 \\
판정 & 접촉·자세·경로 계획·점유·사람 거리 규칙 기반 검출기, 12종 실패 유형 \\
자산화 & 실행별 기록 묶음과 SHA-256 해시, 재실행·비교 실행기 \\
\bottomrule
\end{tabularx}
\vspace{5mm}

{\color{muted} AFS에는 영상 대신 장면 관계와 행동 관측을 텍스트로 제공합니다.}
\end{frame}

% ------------------------------------------------------------------ 6 AFS 입출력
\begin{frame}{능동 실패 탐색(AFS): 무엇을 읽고 무엇을 내놓는가}
\vspace{3mm}
\begin{columns}[T]
\begin{column}{0.48\textwidth}
\textbf{입력: 이전 실행의 행동 관측}
\begin{itemize}
\item 목표 결과·종료 상태
\item 로봇 기능 실행 결과 (경로 계획 실패, 상자 밀기 결과)
\item 물리 측정값 (접촉 시간·접촉력, 물체 이동량)
\item 장면 관계 (상자–목표 ``비어 있음 / 점유'')
\item 이전 제안의 예상 유형과 실제 검출 유형
\end{itemize}
\end{column}
\begin{column}{0.48\textwidth}
\textbf{출력: 검증 가능한 제안}
\begin{itemize}
\item 근거로 삼은 이전 실행
\item 변경할 변수와 값 (1--2개)
\item 실패 가설과 대안 설명
\item 예상 실패 유형 (12종 중)
\item 반증 조건과 측정 계획
\end{itemize}
\end{column}
\end{columns}
\vspace{6mm}
\begin{center}
\colorbox{afs!10}{\parbox{0.92\textwidth}{\centering
규칙 기반 검증기가 범위·근거·남은 시도 횟수를 검사한 뒤 실행 \;→\; 독립된 평가기가 판정\\
{\small\color{muted}예상은 판정에 쓰지 않으므로, LLM 가설의 적중 여부를 사후에 확인할 수 있다}}}
\end{center}
\end{frame}

% ------------------------------------------------------------------ 7 실험 설정
\begin{frame}{두 종류의 테스트 대상}
\vspace{3mm}
\small
\begin{tabularx}{\textwidth}{@{}p{28mm}YY@{}}
\toprule
& 실험 1: 로봇팔 pick-and-place & 실험 2: 휴머노이드 이동 \\
\midrule
로봇 (몸체) & Franka Emika Panda: 7자유도 로봇팔 + 집게형 그리퍼 & Unitree G1: 키 약 1.3\,m 휴머노이드 \\
정책 (두뇌) & GR00T N1.7 (NVIDIA VLA): 카메라 영상·명령 → 관절 동작 & LLM Agent (GPT-6 Luna)가 이동·밀기 기능 선택 + GR00T 보행 정책 \\
작업 & 그릇을 집어 접시에 놓기 & 상자·장애물·횡단하는 사람이 있는 통로에서 목표 도달 \\
환경 변경 & 10개 변수 (그릇·접시·장애물) & 21개 변수 (상자·바닥 물성, 통로 폭, 장애물, 사람) \\
시도 횟수 & 3세트 × 2회 = 방법당 6회 & 3세트 × 6회 = 방법당 18회 (공통 시작 장면은 평가 전 실행, 제외) \\
공통 시작 장면 & 세트별 처음 2회, 같은 무작위 장면 & 세트별 처음 1회, 같은 기본 장면 \\
탐색 LLM & GPT-6 Luna Pro & GPT-6 Luna \\
\bottomrule
\end{tabularx}
\end{frame}

% ------------------------------------------------------------------ 8 결과
\begin{frame}{결과: AFS가 3배 이상 많은 실패 사례를 찾았다}{같은 시도 횟수에서 Random 테스트와 비교}
\begin{columns}[c]
\begin{column}{0.56\textwidth}
\begin{tikzpicture}
\begin{axis}[
  ybar, width=\linewidth, height=62mm, bar width=13pt,
  ymin=0, ymax=118, ylabel={실패 발견율 (\%)},
  symbolic x coords={A,B}, xtick=data,
  xticklabels={{로봇팔},{휴머노이드}},
  xticklabel style={font=\footnotesize,align=center}, tick label style={font=\footnotesize},
  nodes near coords, nodes near coords style={font=\footnotesize},
  legend style={font=\footnotesize,at={(0.5,1.02)},anchor=south,draw=none,legend columns=3},
  axis lines*=left, enlarge x limits=0.18]
\addplot[fill=afs,draw=afs] coordinates {(A,100) (B,33.3)};
\addplot[fill=rnd,draw=rnd] coordinates {(A,16.7) (B,11.1)};
\draw[failc,dashed,thick] (rel axis cs:0,0.254) -- (rel axis cs:1,0.254);
\addlegendimage{legend image code/.code={\draw[failc,dashed,thick] (0cm,0.08cm) -- (0.5cm,0.08cm);}}
\legend{AFS, Random 테스트, 목표 30\%}
\end{axis}
\end{tikzpicture}
\end{column}
\begin{column}{0.4\textwidth}
\textbf{로봇팔} (3세트 × 2회 = 6회)\\
AFS \bignum{100\%}{afs}\\ Random 테스트 16.7\%\\
{\small\color{muted}+500\%, 6배 · 3건은 같은 장면 재제안}\par\vspace{4mm}
\textbf{휴머노이드} (3세트 × 6회 = 18회)\\
AFS \bignum{33.3\%}{afs}\\ Random 테스트 11.1\%\\
{\small\color{muted}+200\%, 3배 · 세 세트 모두 2건씩 (Random 테스트 0/0/2)}\par
\vspace{2mm}
{\scriptsize\color{muted}공통 시작 장면은 평가에서 제외}
\end{column}
\end{columns}
\end{frame}

% ------------------------------------------------------------------ 9 예상 적중
\begin{frame}{LLM 가설은 얼마나 맞았나}
\vspace{3mm}
\begin{columns}[T]
\begin{column}{0.47\textwidth}
\textbf{로봇팔}
\begin{itemize}
\item 6건 모두 ``이송 경로의 장애물이 접근을 막을 것''
\item 6건 모두 로봇–장애물 접촉(10.7--13.7초)을 동반한 FAIL
\end{itemize}
\vspace{3mm}
\textbf{휴머노이드} (제안 18건)
\begin{itemize}
\item 목표 위치 점유 예상 3건 → \textcolor{passc}{\textbf{3건 모두 적중}}
\item 충돌·이동 정체 예상 9건 → 모두 PASS
\end{itemize}
\end{column}
\begin{column}{0.47\textwidth}
\textbf{발견한 취약점}
\begin{itemize}
\item 두 방법의 FAIL 8건이 모두 \textbf{상자 위치} 또는 \textbf{통로 폭} 변경
\item 장애물 형상·바닥 마찰은 취약점이 아니었음
\end{itemize}
\vspace{4mm}
\colorbox{afs!10}{\parbox{0.95\linewidth}{
관측된 실패 메커니즘(차단, 점유)을 겨냥할 때 잘 맞고,\\
아직 관측되지 않은 유형을 겨냥할 때는 빗나감}}
\end{column}
\end{columns}
\end{frame}

% ------------------------------------------------------------------ 10 경계
\begin{frame}{PASS/FAIL 경계: 어디서부터 실패하는가}{로봇팔 통제 실험 11개 장면, 한 번에 한 조건만 변경}
\centering
\resizebox{0.95\textwidth}{!}{\input{build/fig_boundary.tex}}
\vspace{3mm}

\small 장애물 높이 \textbf{12--14\,cm}, 접시 이동 \textbf{37.5--50\%} 사이에서 결과가 바뀐다 → 배포 전 우선 점검 조건
\end{frame}

% ------------------------------------------------------------------ 11 실패 유형
\begin{frame}{실패 유형 다양성: 6종 중 4종 발견 + 도달 불가 부분 관측}
\vspace{3mm}
\small
\begin{tabularx}{\textwidth}{@{}p{30mm}cp{28mm}Y@{}}
\toprule
유형 & 발견 & 실험 & 관측 근거 \\
\midrule
충돌 & \textcolor{passc}{\textbf{O}} & 로봇팔 & 로봇–장애물 접촉 10.7--13.7초, 그릇 미이동 \\
장애물 간섭 & \textcolor{passc}{\textbf{O}} & 휴머노이드 & 장애물 부근 경로 계획 반복 실패, 정체 \\
목표 위치 점유 & \textcolor{passc}{\textbf{O}} & 휴머노이드 & 상자가 목표 영역 점유, 목표 0.85\,m 앞 정체 \\
작업자 안전 위험 & \textcolor{passc}{\textbf{O}} & 휴머노이드 & 사람 1\,m 이내 0.5초 이상 노출, AFS 11/18회 \\
도달 불가 & 부분 & 로봇팔·휴머노이드 & 접시 8.7\,cm 앞 미완료, 최종 목표점 경로 없음 반복 \\
인식 오류 & -- & -- & 후속 실험 (이번엔 실제 위치 값 제공) \\
\bottomrule
\end{tabularx}
\vspace{4mm}

{\color{muted} 12종 체계로 확장 · 확장 유형인 넘어짐은 창고 배치에서 관측}
\end{frame}

% ------------------------------------------------------------------ 12 목표 달성
\begin{frame}{목표 달성 현황}
\vspace{3mm}
\small
\begin{tabularx}{\textwidth}{@{}p{46mm}p{16mm}Y@{}}
\toprule
목표 & 결과 & 근거 \\
\midrule
실패 발견율 30\% 이상 & \textbf{달성} & 로봇팔 100\% (6/6); 휴머노이드 33.3\% (6/18) \\
Random 테스트 대비 20\% 이상 향상 & \textbf{달성} & 로봇팔 +500\%, 휴머노이드 +200\% \\
실패 유형 6종 중 4종 & \textbf{달성} & 충돌, 장애물 간섭, 목표 위치 점유, 작업자 안전 위험 \\
3D 장면 구조화·후보 생성 & 구현 & 10개·21개 장면 변수, 장면 그래프·MJCF·이동 지도 생성 \\
행동 기반 취약 조건 탐색 & 구현 & 관측 → 가설 → 제안 → FAIL 연결 확인 \\
회귀 테스트 자산화 & 구현 & 재실행·비교 실행기, 로봇 LLM 교체 후 결과 유지 확인 \\
\bottomrule
\end{tabularx}
\vspace{4mm}

\textbf{앞으로}\; 현장 디지털 트윈에서 탐색 → 실제 로봇으로 재현 · 실제 인식 결과로 시험
\end{frame}

% ------------------------------------------------------------------ 13 마무리
{\setbeamertemplate{footline}{}
\begin{frame}
\vspace*{12mm}
{\fontsize{24}{30}\selectfont\bfseries 실패를 기다리지 않고, 찾아서 남긴다\par}
\vspace{8mm}
\begin{itemize}
\item 행동 기록을 읽는 LLM이 다음 시험을 설계 \;→\; Random 테스트보다 3배 이상 많은 실패 사례 발견
\item PASS/FAIL 경계를 정량적으로 제시 \;→\; 배포 전 점검 우선순위
\item 발견한 실패는 기록과 함께 저장 \;→\; 정책·배치가 바뀔 때 회귀 테스트로 재실행
\end{itemize}
\end{frame}}

\end{document}
